\section{Introduction and scope}
\label{sec:introduction}

The sequence represented here by TPC-7, TPC-12, and TPC-13 separated
three questions that are often compressed into one.  A complete
multiplicative aggregate can be evaluated
at a fixed shift, but may hide the factor scale.  A Mellin coordinate can
resolve primitive ratio rays, but a ratio alone collapses every radial
multiple.  Adding the GCD coordinate restores the full factor label, but
exact inversion merely reveals the coefficient: it does not estimate the
fixed-shift prime pair sitting at the endpoint.  See
\citet{WangTPC7,WangTPC12,WangTPC13}.

This paper takes the next step in two different quantifier regimes.  First,
we keep one exact radial shell and retain a prescribed fixed shift.  Second,
we keep an arbitrary signed label mask but average the shift over a growing
window.  The distinction is essential.

All asymptotics are as $X\to\infty$.  The functions $W,G$ and the
fixed shift $h$ in the coprime-shell theorem are held fixed; uniform
parameters are identified explicitly.

Put
\begin{equation}
 u(n)=\Lambda(n)-1,
 \qquad
 \eta(t)=\Lambda(t)-1,
 \label{eq:source-target-residuals}
\end{equation}
and fix nonzero real functions
$W\in C_c^1((0,\infty))$ and $G\in C_c^1(\R)$, with
\begin{equation}
 \supp W\subset[\alpha,\beta]
 \subset(0,\infty).
 \label{eq:W-support}
\end{equation}
\begin{equation}
 I_0(W)=\int_0^\infty W(v)\,dv,
 \qquad
 I_0(G)=\int_{\R}G(v)\,dv.
 \label{eq:I0-definitions}
\end{equation}
Every factor pair has the unique representation
\begin{equation}
 n=ak,\qquad m=bk,\qquad (a,b)=1,
 \label{eq:gcd-coordinates}
\end{equation}
and the product center is
\begin{equation}
 r=mn=abk^2.
 \label{eq:product-center}
\end{equation}

\subsection{The fixed-shift radial shell}

The shell $k=1$ consists of the coprime factor pairs $(a,b)$.  At source
cutoff $D=1$, its signed residual is
\begin{equation}
 \cK_h(X)=
 \sum_{\substack{a>1,\ b\ge1\\(a,b)=1}}
 u(a)\eta(ab+h)W(ab/X).
 \label{eq:intro-k-shell}
\end{equation}
For a fixed nonzero integer $h$, define
\begin{equation}
 \beta_h=
 \prod_{p\mid h}
 \left(1-\frac{p}{p^2-p+1}\right),
 \label{eq:beta-h}
\end{equation}
where the product is over the prime divisors of $|h|$ and the empty product
is $1$.  Our first main theorem is
\begin{equation}
 \cK_h(X)=
 \left(\frac1{\zeta(2)}-\beta_h\right)
 XI_0(W)\log X+O_{h,W}(X).
 \label{eq:intro-fixed-shell-law}
\end{equation}
The proof is not a prime-pair estimate.  Grouping by $r=ab$ turns the source
coefficient into
\begin{equation}
 \log\rad(r)-2^{\omega(r)}+1.
 \label{eq:intro-unitary-coefficient}
\end{equation}
The term $2^{\omega(r)}$ is the unitary divisor function.  The fixed-shift
prime target against this function is evaluated by the published
unitary-divisor Titchmarsh theorem of \citet{BiaoWang2026}.  The remaining
terms are elementary at the required precision.

For $h=2$, \eqref{eq:beta-h} gives $\beta_2=1/3$ and the coefficient in
\eqref{eq:intro-fixed-shell-law} is
\begin{equation}
 \frac6{\pi^2}-\frac13>0.
 \label{eq:h-two-shell-coefficient}
\end{equation}
This is a fixed-$2$ theorem, but not a twin-prime theorem.  The shell retains
$k=1$ and sums every primitive orientation $b$; the twin-prime endpoint is
the single orientation $b=1$.

\subsection{The complete dual mask under a short-shift mean}

In the short-shift part we impose the geometric condition that all
translated targets occurring on
\eqref{eq:W-support} remain in a fixed positive dyadic multiple of $X$.
For example, it is enough that
\begin{equation}
 |h_0|+H\max_{v\in\supp G}|v|\leq\frac{\alpha X}{2}.
 \label{eq:translation-geometric-range}
\end{equation}
The one-sided version with $G$ supported in $(0,\infty)$ allows the same
upper range as TPC-12 after changing the fixed dyadic constants.

For an arbitrary shift-independent mask
\begin{equation}
 |\omega_{X,D}(a,b,k)|\leq1,
 \label{eq:bounded-mask}
\end{equation}
define
\begin{equation}
 \cZ_{h,D}^{\omega}(X)=
 \sum_{\substack{a,b,k\geq1,\ (a,b)=1\\ak>D}}
 \omega_{X,D}(a,b,k)u(ak)
 \eta(abk^2+h)W(abk^2/X).
 \label{eq:masked-residual}
\end{equation}
The mask may depend on $X,D,H$ and may be discontinuous in every label, but
it may not depend on the sampled shift $h$.

For fixed $0<\eps<13/15$, uniformly in
\begin{equation}
 X^{2/15+\eps}\leq H\leq X
 \label{eq:short-shift-range}
\end{equation}
subject to \eqref{eq:translation-geometric-range}, we prove that, for every
$A>0$,
\begin{equation}
 \sup_{\|\omega\|_\infty\leq1}
 \left|
 \frac1H\sum_{q\in\Z}G(q/H)
 \cZ_{h_0+q,D}^{\omega}(X)
 \right|
 \ll_{A,\eps,W,G}X(\log X)^{-A}.
 \label{eq:intro-mask-theorem}
\end{equation}
This is stronger than uniformity in the two character phases
$(a/b)^{it}k^{is}$.  It closes that phase family to the entire bounded
$\ell^\infty$ coefficient-mask ball.  In particular it includes the exact endpoint mask
$\one_{b=k=1}$, a single radial layer $\one_{k=k_0}$, one primitive ray,
a dyadic source layer, and arbitrary signed unions of these sets.

The analytic input is Corollary~1.4 of \citet{GuthMaynard2026}, which gives
the prime number theorem in almost all intervals of length
$X^{2/15+\eps}$.  Our contribution is an $\ell^2$ dual transfer.  The
translated residual vector has arbitrarily small logarithmic norm, while
every bounded GCD--Mellin mask has a divisor-controlled $\ell^2$ norm.
This is not a new short-interval theorem.

\subsection{Endpoint, hard packets, and the information gate}

The endpoint mask yields
\begin{equation}
 \frac1H\sum_qG(q/H)
 \sum_n u(n)\eta(n+h_0+q)W(n/X)
 \ll_A X(\log X)^{-A}.
 \label{eq:intro-endpoint-covariance}
\end{equation}
After adding back the one-prime terms, its raw form has main term
$XI_0(W)I_0(G)$.  This reaches the endpoint before taking the mean, but the
mean still samples $\asymp H$ shifts.  It gives no information about the
coordinate $h=2$.

There is also a legal target-side Vaughan formulation.  One first uses the
factorized target $t=\ell k$ as the outer center; only then may one apply an
arbitrary mask $\omega_{\ell,k}$.  The coefficient is independent of the
subsequent translation, so \eqref{eq:intro-mask-theorem}'s abstract
$\ell^2$ mechanism gives a short-shift averaged model for every masked hard
packet.  Masking a factorization of $r+h$ while retaining $r$ as center
would instead make the mask $h$-dependent and is not covered.

Finally we give two exact boundary theorems.  First, nonlinear recovery from
a specified local pushforward has minimax radius equal to an
$\ell^\infty$ row-space distance.  Second, the complete divisor signature
up to $R$ distinguishes primes from composites in $(X,2X]$ once
$R\geq\sqrt{2X}$, whereas below
$(1-\eps)\sqrt{2X}$ a prime and a squarefree semiprime have identical zero
signatures.  A balanced bilinear factor row separates this null pair exactly.
That is a test that the row contains new information, not a proof that its
arithmetic sum is small.

\subsection{Claim ledger and relation to earlier layers}

\begin{center}
\small
\begin{tabular}{
 >{\raggedright\arraybackslash}p{0.18\textwidth}
 >{\raggedright\arraybackslash}p{0.25\textwidth}
 >{\raggedright\arraybackslash}p{0.22\textwidth}
 >{\raggedright\arraybackslash}p{0.18\textwidth}}
\toprule
Result & Label retained & Quantifier & Missing conclusion\\
\midrule
TPC-7 & complete large-divisor defect & almost all $h$, $H\geq X^{8/33+\eps}$ & prescribed $h$ and individual layers\\
TPC-13 & every $(a,b,k)$ through character phases & signed first moment and squared energy & arbitrary bounded masks\\
Fixed shell here & exact $k=1$ shell & every fixed $h\ne0$ & one orientation $b=1$\\
Mask theorem here & arbitrary $(a,b,k)$ mask & signed mean, $H\geq X^{2/15+\eps}$ & any selected shift\\
\bottomrule
\end{tabular}
\end{center}

The constructed GCD--Mellin and translation actions below are coordinate
representations of finite coefficient arrays.  We claim no intrinsic prime,
logistic, transfer-operator, or zeta-zero spectrum.  We also prove no
fixed-shift prime-pair asymptotic, twin-prime lower bound, new level of
distribution, or violation of the classical sieve parity barrier.
